给定, 零假设为 $H_0:\boldsymbol{\mu}'=[7,11]$, 试计算 $T^2$.

(b)\ 确定 (a) 中 $T^2$ 的分布.

(c)\ 利用 (a) 与 (b), 在 $\alpha=0.05$ 水平下检验 $H_0$. 你得出什么结论?

\noindent 5.2\quad 利用例 5.1 的数据, 验证: 若将每个观测值 $\mathbf{x}_j$, $j=1,2,3$, 分别用 $\mathbf{C}\mathbf{x}_j$ 替代, 则 $T^2$ 保持不变. 此处
\[
\mathbf{C}=\begin{bmatrix}1&-1\\1&1\end{bmatrix}
\]
注意: 观测值
\[
\mathbf{C}\mathbf{x}_j=\begin{bmatrix}x_{j1}-x_{j2}\\x_{j1}+x_{j2}\end{bmatrix}
\]
产生的数据矩阵为:
\[
\begin{bmatrix}
(6-9)&(10-6)&(8-3)\\
(6+9)&(10+6)&(8+3)
\end{bmatrix}'
\]

\noindent 5.3\quad (a) 对练习 5.1 中的数据, 利用\hyperref[eq:5-15]{式(5-15)}计算 $T^2$.

(b)\ 用练习 5.1 的数据, 计算\hyperref[eq:5-13]{式(5-13)}中的 $\Lambda$, 并计算威尔克斯 $\Lambda$ 值.

\noindent 5.4\quad 利用\hyperref[tab:5.1]{表 5.1} 的汗液数据(见例 5.2)

(a)\ 确定 $\boldsymbol{\mu}$ 的 $90\%$ 置信椭球的主轴, 并计算其主轴的长度.

(b)\ 对排汗量、钠含量与钾含量的数据, 分别作 Q-Q 图, 并对每一对变量的观测值作散布图. 对这种情况, 多元正态假定是否合理? 试讨论.

\noindent 5.5\quad 对例 5.3 中的微波炉辐射数据给出的 $\bar{\mathbf{x}}$, $\mathbf{S}$ 及 $\mathbf{S}^{-1}$, 在显著性水平 $\alpha=0.05$ 下检验零假设 $H_0:\boldsymbol{\mu}'=[0.55,0.60]$, 所得结果是否与\hyperref[fig:5.1]{图 5.1} 中 $\boldsymbol{\mu}$ 的 $95\%$ 置信椭圆相一致? 请作出解释.

\noindent 5.6\quad 对 $m=3$ 验证庞弗罗尼不等\hyperref[eq:5-28]{式(5-28)}. 提示: 考虑三个事件 $C_1$, $C_2$ 与 $C_3$ 的韦恩图也许有帮助.

\noindent 5.7\quad 对\hyperref[tab:5.1]{表 5.1} 的汗液数据(见例 5.2), 利用结论 5.3 计算 $\mu_1$, $\mu_2$ 与 $\mu_3$ 的 $95\%$ 联合 $T^2$ 置信区间, 并用\hyperref[eq:5-29]{式(5-29)}构造其 $95\%$ 庞弗罗尼区间, 并对这两种区间进行比较.

\noindent 5.8\quad 从\hyperref[eq:5-23]{式(5-23)}我们知道, $T^2$ 等于由线性组合 $a'\bar{\mathbf{x}}$ 构造的一元 $t$ 值平方中的最大者, 此处 $a=\mathbf{S}^{-1}(\bar{\mathbf{x}}-\boldsymbol{\mu})$. 利用例 5.3 中的结论及练习 5.5 中的 $H_0$, 对变换后的微波炉辐射数据计算 $a$ 值. 验证: 由此 $a$ 值计算出的 $t^2$ 值等于练习 5.5 中的 $T^2$ 值.

\noindent 5.9\quad 为了维持大灰熊种群健康, 阿拉斯加州渔业和狩猎部的自然学家哈利·罗伯茨对大灰熊进行了研究. 在观测 61 只大灰熊的基础上, 他得到了如下的汇总统计数据(在练习 8.23 里也可看到):

\begin{center}
\begin{tabular}{lcccccc}
\toprule
变量 & 重量(kg) & 体长(cm) & 颈(cm) & 腰围(cm) & 头长(cm) & 头宽(cm)\\
\midrule
样本均值 $\bar{x}$ & 95.52 & 164.38 & 55.69 & 93.39 & 17.98 & 31.13\\
\bottomrule
\end{tabular}
\end{center}

协方差矩阵为:
\[
\mathbf{S}=\begin{bmatrix}
3266.46&1343.97&731.54&1175.50&162.68&238.37\\
1343.97&721.91&324.25&537.35&80.17&117.73\\
731.54&324.25&179.28&281.17&39.15&56.80\\
1175.50&537.35&281.17&474.98&63.73&94.85\\
162.68&80.17&39.15&63.73&9.95&13.88\\
238.37&117.73&56.80&94.85&13.88&21.26
\end{bmatrix}
\]
